\section{Hybrid System Architecture}

The architecture is a one-way pipeline:

\begin{enumerate}
  \item finite-fault source and terrain preprocessing;
  \item Regional2D finite-volume propagation;
  \item extraction of a nearshore coupling interface;
  \item mapping to Local3D OpenFOAM replay boundary data;
  \item local no-defence and rigid-barrier hydrodynamic runs;
  \item evidence and provenance capture.
\end{enumerate}

\begin{table}[H]
\centering
\caption{Traceability classification for major mathematical components.}
\label{tab:traceability-classification}
\begin{tabularx}{\textwidth}{p{0.24\textwidth}p{0.22\textwidth}X}
\toprule
Component & Implementation route & Evidence and limitation \\
\midrule
Regional NLSWE & Repository-native & Research authority in RES-MOD/RES-NUM; G6 prefix equivalence accepted. \\
Earthquake preprocessing & Repository-owned preprocessing & Approved finite-fault input; source alternatives remain calibration sensitivity candidates. \\
Corridor/geospatial preprocessing & Repository-owned & EPSG:32654 corridor and conditioned terrain; higher-resolution products remain post-G6. \\
2D to 3D coupling & Repository-native replay/mapping & Fixed accepted source window 245--545 s mapped to 0--300 s. \\
Local3D URANS--VOF & Adopted OpenFOAM Foundation 11 backend & Repository defines dictionaries and evidence; OpenFOAM internals are adopted, not repository-native. \\
Wall functions & OpenFOAM-adopted with project policy & Continuous Spalding baseline; formal wall/mesh convergence remains post-G6. \\
Boundary treatment & Project-defined OpenFOAM configuration & \texttt{open\_ocean\_damped} policy accepted by reflection benchmark and real replay. \\
\bottomrule
\end{tabularx}
\end{table}
